% =========================================================================
% Abstract (qualitative only: no digits, concentrations, yields, or zone sizes)
% =========================================================================

\clearpage
\thispagestyle{frontmatterstyle}
\begin{center}
{\Large \bfseries ABSTRACT\par}
\vspace{0.25in}
\end{center}

\noindent Antimicrobial resistance among Gram-negative bacteria remains a clinical problem. The outer membrane limits the entry of many drugs. Pumps and enzymes then reduce the effect of those that do enter. \taxa{Escherichia coli}, \taxa{Pseudomonas aeruginosa}, and \taxa{Klebsiella pneumoniae} cause urinary, wound, lung, and bloodstream infection. Plant extracts are therefore one source of antibacterial compounds.

\noindent Medicinal plants contain alkaloids, flavonoids, phenolics, tannins, saponins, terpenoids, and related glycosides. These classes can disturb membranes, bind proteins, and interfere with adhesion. Water is the solvent used in many household decoctions. An aqueous extract is therefore closer to traditional use than a dichloromethane or methanol extract, even if water recovers a different set of constituents.

\noindent \taxa{Dipsacus inermis} Wall., the Himalayan teasel, belongs to Caprifoliaceae. It grows in the Kashmir Himalaya and is known locally as wopal haakh. Folklore records use of the leaves for inflammation, cough, and postnatal bathing. Earlier laboratory work on this species has focused on organic-solvent extracts and on anti-inflammatory models. The antibacterial behaviour of a simple aqueous Soxhlet extract, read only as a zone of inhibition on seeded agar, has received less attention. \taxa{Klebsiella pneumoniae} has not been a regular test organism in that literature.

\noindent The present study was carried out at the Advanced Research Laboratory, Department of Zoology, University of Kashmir. Aerial material was collected from the Kashmir University Botanical Garden, authenticated at the Centre for Biodiversity and Taxonomy, shade-dried, powdered, and extracted with distilled water in a Soxhlet apparatus. The extract was concentrated and reconstituted at a lower working strength and at an elevated working strength. Qualitative tube tests recorded the main metabolite classes. Antibacterial activity was then scored only as the zone of inhibition against the three named pathogens. Distilled water was the negative control. A gentamicin GEN 5 paper disc was the positive control.

\noindent Phytochemical tests gave positive reactions for flavonoids, terpenoids, quinones, alkaloids, tannins, saponins, phenols, carbohydrates, and phytosterols. Saponin tubes held a persistent froth. Distilled water produced no zone of clearance on any plate, which indicates that the solvent itself did not suppress growth. Both extract strengths produced visible zones. The lower strength gave a pronounced zone against \taxa{Pseudomonas aeruginosa}. The elevated strength enlarged the zone against \taxa{Escherichia coli} and, more modestly, against \taxa{Klebsiella pneumoniae}, but did not enlarge the zone against \taxa{Pseudomonas aeruginosa}. The gentamicin disc produced a wide zone on every plate and remained the strongest comparator.

\noindent The pattern is consistent with polar phenolics, tannins, flavonoids, and saponins recovered by water, acting against organisms that still possess an outer-membrane barrier. It does not support a uniform rise in activity with extract strength, and it does not replace a minimum inhibitory concentration assay. Within those limits, the aqueous extract of \taxa{Dipsacus inermis} showed \invitro{} antibacterial activity against the three Gram-negative isolates tested here.

\vspace{0.2in}
\noindent \textbf{Key words:} \taxa{Dipsacus inermis}, aqueous extract, phytochemical screening, disc diffusion, zone of inhibition, gentamicin, Gram-negative bacteria, Kashmir.
